\section{Time Hierarchy --- Separation}
%%% generated by the blueprint pipeline from TCSlib.Complexity.TimeHierarchy.Separation
\subsection*{Time Hierarchy --- Separation: Introduction}

Declarations of the Lean module \texttt{TCSlib.Complexity.TimeHierarchy.Separation}.

\subsection*{Time Hierarchy --- Separation: Results}

\begin{definition}[Binary Output Machine]
\lean{Complexity.twoPowTM}
\label{Complexity.twoPowTM}
\leanok
\uses{Turing.FinTM}
A finite Turing machine over the binary alphabet with no work tapes and one state is
defined as follows: when it scans an input symbol, it moves the input head right,
outputs $0$, and remains in its state; when it scans a blank, it outputs $1$ and halts.
\end{definition}

\begin{lemma}[Run of the unary-output machine]
\lean{Complexity.twoPowTM_run}
\label{Complexity.twoPowTM_run}
\leanok
\uses{Complexity.twoPowTM, Turing.Action.apply, Turing.Cfg.inputSymbol, Turing.MultiTapeTM.initCfg, Turing.MultiTapeTM.runFrom, Turing.MultiTapeTM.runFrom_succ_eq_step', Turing.MultiTapeTM.step, Turing.inputSymbolInner, Turing.moveInputPos_pos_of_ne_right}
\difficulty{6}
For every binary word $x$ and every $i\le |x|$, after $i$ steps the specified machine
started on $x$ is in its unique state, its input head is at position $i+1$, and its
output consists of $i$ zeros.
\end{lemma}

\begin{lemma}[Binary digits of a power of two]
\lean{Complexity.bits_two_pow}
\label{Complexity.bits_two_pow}
\leanok
\difficulty{4}
For every natural number $n$, the binary digit list of $2^n$, written from least
significant to most significant digit, consists of $n$ zeros followed by a single $1$.
\end{lemma}

\begin{theorem}[Exponential time constructibility]
\lean{Complexity.timeConstructible_two_pow}
\label{Complexity.timeConstructible_two_pow}
\leanok
\uses{Complexity.TimeConstructible, Complexity.bits_two_pow, Complexity.twoPowTM, Complexity.twoPowTM_run, Turing.Action.apply, Turing.Cfg.inputSymbol, Turing.FinTM.ComputesInTime, Turing.FinTM.ComputesInTime.mono, Turing.FinTM.computesInTime_iff, Turing.MultiTapeTM.initCfg, Turing.MultiTapeTM.runFrom, Turing.MultiTapeTM.runFrom_succ_eq_step', Turing.MultiTapeTM.step}
\difficulty{5}
The function $T(n)=2^n$ is time constructible: for every $n$, $n\le 2^n$, and some
positive integer constant $a$ and finite binary Turing machine compute, on every binary
input $x$, the little-endian binary representation of $2^{|x|}$ within $a(2^{|x|}+1)$
steps.
\end{theorem}

\begin{theorem}[A polynomial is eventually bounded by an exponential]
\lean{Complexity.eventually_poly_le_two_pow}
\label{Complexity.eventually_poly_le_two_pow}
\leanok
\difficulty{6}
For all natural numbers $A$ and $K$, there is a natural number $N$ such that for every
$n\ge N$, $A(n+1)^K\le 2^n$.
\end{theorem}

\begin{theorem}[Polynomial Square Eventually Bounded by an Exponential]
\lean{Complexity.eventually_poly_sq_le_two_pow}
\label{Complexity.eventually_poly_sq_le_two_pow}
\leanok
\uses{Complexity.eventually_poly_le_two_pow}
\difficulty{5}
For natural numbers $k$ and $A$, there is a natural number $N$ such that for every $n
\ge N$, $A(n^k+1+n+1)^2 \le 2^n$.
\end{theorem}

\begin{theorem}[Strict polynomial versus exponential time inclusion]
\lean{Complexity.dtime_poly_ssubset_dtime_two_pow}
\label{Complexity.dtime_poly_ssubset_dtime_two_pow}
\leanok
\uses{Complexity.DTIME, Complexity.eventually_poly_sq_le_two_pow, Complexity.timeConstructible_two_pow, Complexity.time_hierarchy_of_pos}
\difficulty{1}
For every natural number $k$, $\mathrm{DTIME}(n\mapsto n^k+1)$ is a proper subset of
$\mathrm{DTIME}(n\mapsto 2^n)$.
\end{theorem}

\begin{theorem}[Strict containment in exponential time]
\lean{Complexity.P_ssubset_EXP}
\label{Complexity.P_ssubset_EXP}
\leanok
\uses{Complexity.EXP, Complexity.P, Complexity.P_subset_EXP, Complexity.TimeHierarchy.diagLang, Complexity.TimeHierarchy.diagLang_mem_DTIME, Complexity.TimeHierarchy.diagLang_not_mem_DTIME, Complexity.eventually_poly_sq_le_two_pow, Complexity.timeConstructible_two_pow, Turing.FinTM.ComputesInTime.mono}
\difficulty{5}
The class $\mathsf{P}$ of polynomial-time decidable binary languages is a proper subset
of the class $\mathrm{EXP}$, consisting of languages decidable in time $2^{n^c}$ up to a
constant factor for some natural number $c$.
\end{theorem}

\begin{theorem}[Separation of polynomial and exponential time]
\lean{Complexity.P_ne_EXP}
\label{Complexity.P_ne_EXP}
\leanok
\uses{Complexity.EXP, Complexity.P, Complexity.P_ssubset_EXP}
\difficulty{1}
The class $\mathsf{P}$ of languages decidable in polynomial time is not equal to the
class $\mathrm{EXP}$ of languages decidable in time $2^{n^c}$ up to a constant factor,
for some natural number $c$.
\end{theorem}
